\documentclass[11pt]{article}
\input{preamble.tex}

\title{\huge{Oblig 2}}
\author{\LARGE{Thobias Høivik}}
\date{}

\begin{document}
\maketitle

\newpage
\begin{prob*}[1]
    La $G$ of $H$ være endelige grupper, og la $\phi : G \to H$
    være en gruppehomomorfi.

    \begin{enumerate}[label=(\alph*), align=left]
        \item Bruk Lagranges teorem og fundamentalteoremet for 
            gruppehomomorfier til å vise at $|G|/|\ker \phi|$ deler 
            $|H|$.
        \item Anta nå at $\gcd(|G|,|H|) = 1$. Vis at da er 
            $\phi$ den trivielle homomorfien, det vil si 
            $\phi(g) = e$ for alle $g \in G$.
    \end{enumerate}
\end{prob*}
\begin{proof}[Bevis av (a)]
    Ved Lagranges teorem anvendt på undergruppen $\ker\phi \le G$ får vi
    \[
        |G| = |\ker\phi|\,[G:\ker\phi].
    \]
    Dermed er
    \[
        \frac{|G|}{|\ker\phi|} = [G:\ker\phi] = |G/\ker\phi|.
    \]

    Ved fundamentalteoremet for gruppehomomorfier har vi
    \[
        G/\ker\phi \cong \operatorname{Im}(\phi),
    \]
    og derfor
    \[
        |G/\ker\phi| = |\operatorname{Im}(\phi)|.
    \]
    Vi får altså
    \[
        \frac{|G|}{|\ker\phi|} = |\operatorname{Im}(\phi)|.
    \]

    Siden $\operatorname{Im}(\phi)$ er en undergruppe av $H$, 
    følger det av
    Lagranges teorem at
    \[
        |\operatorname{Im}(\phi)| \bigm| |H|.
    \]
    Dermed har vi vist at
    \[
        \frac{|G|}{|\ker\phi|} \biggm| |H|.
    \]
\end{proof}
\begin{proof}[Bevis av (b)]
    Vi vet at 
    \[
        \frac{|G|}{|\ker \phi|} \biggm| |H|,
    \]
    men også at 
    \[
        |G| = |\ker \phi| \cdot \frac{|G|}{|\ker\phi|},
    \]
    så $|G|/|\ker \phi|$ deler $|G|$. Altså $|G|/|\ker \phi|$ deler 
    både $|G|$ og $|H|$ med $\gcd(|G|,|H|) = 1$ så 
    $|G|/|\ker \phi| = 1$ så $|\operatorname{Im}(\phi)| = 1$ altså 
    \[
        \operatorname{Im}(\phi) = \{e\}
    \]
\end{proof}

\newpage 
\begin{prob}[2]
    La $R$ og $S$ være ringer med identitetselementer $1_R \in R$ 
    og $1_S \in S$, slik at $1_S \neq 0$. La $\phi : R \to S$
    være en ringhomomorfi slik at $\phi(1_R) = 1_S$.

    \begin{enumerate}[label=(\alph*), align=left]
        \item Vis at hvis $r \in R$ er en enhet, så er 
            $\phi(r)$ en enhet i $S$.
        \item Vis at hvis $R$ er en kropp, så er $\phi$ injektiv. 
        \item Vis at den motsatte implikasjonen fra punkt (a) ikke 
            er sann, det vil si: 

            Finn et eksempel på en ringhomomorfi $\phi : R \to S$ 
            med $\phi(1_R) = 1_S \neq 0$ og en $r \in R$
            slik at $\phi(r)$ er en enhet i $S$, men $r$ er ikke
            en enhet i $R$.

            \emph{Hint:} Du kan finne et eksempel med $R = \mathbb Z$.
    \end{enumerate}
\end{prob}
\begin{proof}[Bevis av (a)]
    La $r \in R$ være en enhet, altså det fins 
    $r^{-1} \in R$ slik at 
    \[
        rr^{-1} = r^{-1}r = 1_R
    \]
    Da ser vi at 
    \begin{align*}
        \phi(r)\phi(r^{-1}) &= \phi(rr^{-1}) \\
                            &= \phi(1_R) \\ 
                            &= 1_S
    \end{align*}
    og lignende med $\phi(r^{-1})\phi(r)$.
    Dermed er $\phi(r)$ en enhet i $S$ med invers 
    $\phi(r^{-1})$.
\end{proof}
\begin{proof}[Bevis av (b)]
    $\ker \phi$ er et ideal i $R$. $R$ er en kropp så de enste 
    idealen er $\{0\}$ og $R$. $\phi(1_R) = 1_S \neq 0$ så 
    $1_R \not\in \ker \phi$ altså $\ker \phi \neq R$.
    Vi konkluderer da at $\ker \phi = \{0\}$ dermed injektiv.
\end{proof}
\begin{proof}[Bevis av (c)]
    La $R = \mathbb Z$ og la $S = \mathbb Q$. Definer $\phi : R \to S$ 
    ved 
    \[
        \phi(a) = a
    \]
    Da er $\phi(2) = 2$ som har en multiplikativ invers i $\mathbb Q$, 
    men ikke i $\mathbb Z$.
\end{proof}

\newpage 
\begin{prob}[3]
    Se på produktringen $\mathbb Z \times \mathbb Z$, og la 
    \[
        R = \{(m,n) \in \mathbb Z \times \mathbb Z \bigm| 
        m\equiv n \!\!\!\! \mod 2\}
        \subseteq \mathbb Z \times \mathbb Z.
    \]


    \begin{enumerate}[label=(\alph*), align=left]
        \item Vis at $R$ er en underring av $\mathbb Z \times 
            \mathbb Z$. 
        \item Finn alle $0$-divisorer i $R$.
    \end{enumerate}
\end{prob}
\begin{proof}[Bevis av (a)]
    Merk at $(0,0) \in R \neq \emptyset$ siden $0 \equiv 0$.
    La $(a,b),(c,d) \in R$. Altså $a \equiv b$ og $c \equiv d$.
    Da oppfyller 
    \[
        (a,b)-(c,d) = (a-c,b-d) 
    \]
    $a-c \equiv b-d$. 
    Dermed er $(a-c,b-d)\in R$.

    Videre er
    \[
        (a,b)(c,d) = (ac,bd).
    \]
    Nå har vi
    \[
        ac-bd = c(a-b)+b(c-d).
    \]
    Siden $2 \mid (a-b)$ og $2 \mid (c-d)$, følger det at
    $2 \mid c(a-b)$ og $2 \mid b(c-d)$. Derfor er også
    $2 \mid (ac-bd)$, altså
    \[
        ac \equiv bd \pmod 2.
    \]
    Dermed er $(ac,bd)\in R$.

    Vi har altså vist at $R$ er ikke-tom og lukket under subtraksjon
    og multiplikasjon, så $R$ er en underring av $\mathbb Z \times 
    \mathbb Z$.
\end{proof}
\begin{proof}[Løsning av (b)]
    Vi må finne $(a,b),(c,d) \in R\setminus \{(0,0)\}$ slik at 
    \[
        (a,b)(c,d) = (ac,bd) = (0,0)
    \]
    med $ac,bd \in \mathbb Z$. For $a = 0$ krever vi at 
    $d = 0$ og vice versa for $b$ og $c$, siden ellers ville 
    ene elementet vert $(0,0)$. Siden en av innslagene i tupplen 
    er $0$ og vi krever $m \equiv n$ må det andre være 
    i $2\mathbb Z \setminus \{0\}$. Så $0$-divisorene er alle par 
    \[
        (x,0),(0,y), \ x, y \in 2\mathbb Z \setminus \{0\}. 
    \]
\end{proof}
\newpage 
\begin{prob}[4]
    La $F$ være en kropp. Vi sier at $a \in F$ er et kvadrat 
    hvis det finnes $b$ slik at $a = b^2$. 

    \begin{enumerate}[label=(\alph*), align=left]
        \item Vis at $a \in F$ er et kvadrat hvis og bare hvis 
            polynomet $x^2 - a \in F[x]$ er redusibelt.

            \medskip 
            La nå $p \geq 3$ være et primtal og la $a \in \mathbb Z_p$
            med $a \neq 0$.
        \item Vis at hvis $a$ er et kvadrat, så er 
            $a^{\frac{p-1}{2}} = 1$.
        \item Vis at hvis $a^{\frac{p-1}{2}} = 1$, så er 
            $a$ et kvadrat.
    \end{enumerate}
\end{prob}
\begin{proof}[Bevis av (a)] 
    Vi lar $F$ være en kropp.
        
    ($\Rightarrow$) 

    Anta at $a \in F$ er et kvadrat. Altså det finnes $b \in F$ 
    slik at $a = b^2$.  
    La $f(x) = x^2 - a \in F[x]$. Siden $\deg f = 2$ er 
    $f$ redusibelt hvis og bare hvis det finnes $x \in F$ slik at 
    $f(x) = 0$. Vi observerer at 
    \begin{align*}
        f(b) &= b^2 - a \\ 
             &= a - a \\ 
             &= 0
    \end{align*}
    så $f$ er redusibelt. 

    ($\Leftarrow$) 

    Anta at $f(x) = x^2 - a \in F[x]$ er redusibelt. Siden 
    $\deg f = 2$ er $f$ redusibelt hvis og bare hvis det finnes 
    $x \in F$ slik at $f(x) = 0$. $f$ er redusibelt så det finnes 
    $b \in F$ slik at 
    \begin{align*}
        f(b) &= b^2 - a \\ 
             &= 0 \\ 
        \Rightarrow b^2 &= a
    \end{align*}
\end{proof}
\begin{proof}[Bevis av (b)]
    La $p \geq 3$ være et primtal.
    Anta at $a \in \mathbb Z_p \setminus \{0\}$ er et kvadrat, 
    altså det finnes 
    $b \in \mathbb Z_p$ slik at $a = b^2$. Da har vi  
    \begin{align*}
        a^{\frac{p-1}{2}} &= (b^2)^{\frac{p-1}{2}} \\ 
                          &= b^{p-1}
    \end{align*}
    Via Fermat's teorem: 
    \[
        b^{p-1} \equiv 1 \mod p
    \]
\end{proof}
\begin{proof}[Bevis av (c)]
    Siden $a \neq 0$, er $a \in \mathbb Z_p^\times$. Gruppen
    $\mathbb Z_p^\times$ er syklisk av orden $p-1$, så det finnes
    en generator $g$ slik at
    \[
        \mathbb Z_p^\times = \langle g\rangle.
    \]
    Dermed finnes et heltall $k$ slik at
    \[
        a = g^k.
    \]

    Anta nå at
    \[
        a^{\frac{p-1}{2}} = 1.
    \]
    Da får vi
    \[
        1 = (g^k)^{\frac{p-1}{2}} = g^{k\frac{p-1}{2}}.
    \]
    Siden $g$ har orden $p-1$, følger det at
    \[
        p-1 \mid k\frac{p-1}{2}.
    \]
    Dermed må $2 \mid k$, altså er $k=2m$ for et heltall $m$.

    Da får vi
    \[
        a = g^k = g^{2m} = (g^m)^2.
    \]
    Altså er $a$ et kvadrat i $\mathbb Z_p$.
\end{proof}

\end{document}
